\begin{textsample}{NEG Dim 3 – 2009\_c4 – Score 1.00 – 2009\_c4/enade\_2009\_secr\_exec\_8db19b640f.txt}  \label{ex:f3_neg_164}
Capítulo 4 Percepção sobre a Prova As percepções dos estudantes sobre a prova aplicada no ENADE/2009 na área de SECRETARIADO EXECUTIVO foram mensuradas por meio de nove itens que avaliaram desde o grau de dificuldade da prova até o tempo gasto para concluí -la.
As questões foram analisadas separando -se concluintes e ingressantes, e as percepções sobre a prova foram relacionadas com o desempenho dos estudantes e com a Região Geográfica de localização do curso6.
O desempenho dos estudantes dentro de cada Área/Subárea foi classificado em quartos ( 1/4 ) para os concluintes e para os ingressantes em separado ( ver Figura 1 para o esquema de representação gráfica destas grandezas ).
As notas finais foram colocadas em ordem ascendente e no primeiro quarto foram colocadas as 25 \% notas mais baixas, limitadas superiormente pelo percentil 25 que vem a ser o primeiro quartil.
No segundo quarto, as 25 \% seguintes, limitadas superiormente pelo percentil 50 que vem a ser a mediana.
No terceiro quarto, as 25 \% notas após as encontradas no segundo quarto, ou seja, as primeiras 25 \% acima da mediana, limitadas superiormente pelo percentil 75 que vem a ser o terceiro quartil.
No último quarto foram alocadas as 25 \% maiores notas.
Potencialmente, poderia existir uma relação entre a nota do estudante e a sua percepção sobre a prova.
Calculando as médias das respostas segundo os quartos de notas, seria possível reconhecer a existência ( ou não ) de algum gradiente nas respostas, por exemplo, uma melhor nota correspondendo a uma percepção menos negativa.
Este gradiente seria traduzido, por exemplo, no gráfico 4.1a, pela diminuição das áreas laranja e vermelha ( Difícil e Muito difícil ) com o aumento do desempenho ( indo do primeiro ao último quarto ). 6 Curso, neste relatório, refere -se à unidade de análise para o Conceito ENADE e é caracterizado pela combinação de IES, Área/Subárea ( quando pertinente ) e município de habilitação.
A medida tradicional de “ grau de discriminação ” comparou os resultados dos estudantes do primeiro quarto ( 25 \% piores estudantes ) aos do último quarto ( 25 \% melhores ).
Esta comparação, no gráfico disponibilizado, foi idêntica à comparação da primeira barra ( primeiro quarto ) com a última barra ( último quarto ) de cada grupo ingressante ou concluinte nos gráficos “ segundo quartos de nota final ” ( gráficos com números ímpares neste capítulo ).
Por outro lado, a disponibilização das informações segundo os quartos de nota permitiu verificar o comportamento de dois outros grupos intermediários que na análise somente dos extremos tem que ser suposto monotônico.
A seguir, serão apresentados os principais resultados relativos aos nove itens avaliados. 4.1 Grau de dificuldade da prova 4.1.1 Formação Geral Os Gráficos 4.1a, 4.1b e 4.1c apresentam a distribuição das respostas à questão “ Qual o grau de dificuldade desta prova na parte de Formação Geral? ”, segundo quartos de nota final e grupo de estudantes concluintes e ingressantes que participaram do ENADE/2009 na área de SECRETARIADO EXECUTIVO.
Os Gráficos 4.2a e 4.2b apresentam a distribuição das respostas à questão “ Qual o grau de dificuldade desta prova na parte de Formação Geral? ”, segundo a Região onde o curso está situado para os estudantes concluintes e ingressantes que participaram do ENADE/2009 na área de SECRETARIADO EXECUTIVO.
Entre os concluintes, a maior proporção de estudantes que considerou a parte de Formação Geral da prova como de dificuldade média entre todas as Regiões brasileiras, foram os da Região Norte com 69,4 \%, já entre os ingressantes, foram os estudantes da Região Norte, também, com 67,6 \%.
A Região com as menores proporções desta resposta foram respectivamente para os concluintes e ingressantes, a Centro-Oeste ( 55,5 \% e 57,3 \% ). 4.1.2 Componente Específico Os Gráficos 4.3a, 4.3b e 4.3c apresentam a distribuição das respostas à questão “ Qual o grau de dificuldade desta prova na parte do Componente Específico? ”, segundo quartos de nota final e grupo de estudantes concluintes e ingressantes que participaram do ENADE/2009 na área de SECRETARIADO EXECUTIVO.
Em geral, para as desagregações consideradas, os estudantes classificaram o Componente Específico como mais difícil que a Formação Geral.
No total de concluintes, 24,4 \% declararam que a prova na parte do Componente Específico foi Difícil/Muito difícil, sendo um valor muito maior para os ingressantes, 35,7 \%.
Entre os concluintes, a proporção de indivíduos que classificaram esta parte como “ Muito difícil ” diminui com o desempenho do estudante.
A proporção que declara Difícil/Muito difícil também entre os concluintes diminuiu com o desempenho, indo de 28,8 \% a 22,4 \% no último quarto.
A proporção que declarou a parte Fácil/Muito fácil foi maior entre os concluintes ( 11,2 \% ) do que entre os ingressantes ( 7,5 \% ).
Entre os concluintes a proporção dos que consideraram esta parte como Fácil decresceu com o desempenho.
Entre os concluintes, 64,4 \% classificaram esta parte com um nível médio de dificuldade.
A proporção entre os ingressantes foi menor, 56,8 \%.
Os Gráficos 4.4a e 4.4b apresentam a distribuição das respostas à questão “ Qual o grau de dificuldade desta prova na parte do Componente Específico? ”, segundo a Região onde o curso está situado para os estudantes concluintes e ingressantes que participaram do ENADE/2009 na área de SECRETARIADO EXECUTIVO.
Nenhum concluinte ou ingressante da Região Centro-Oeste e nenhum ingressante da Região Norte declarou que esta parte foi Muito fácil.
Entre os concluintes, a maior proporção de estudantes que considerou a parte do Componente Específico da prova como de dificuldade média entre todas as Regiões brasileiras, foi na Região Norte com 68,1 \% e entre os ingressantes foram os estudantes da Região Nordeste com 64,3 \%.
A Região com as menores proporções desta resposta foram respectivamente para os concluintes e ingressantes, a Centro-Oeste ( 60,3 \% ) e a Sul ( 52,7 \% ). 4.2 Extensão da prova em relação ao tempo total Os Gráficos 4.5a, 4.5b e 4.5c apresentam a distribuição das respostas à questão sobre a extensão da prova em relação ao tempo total, segundo quartos de nota final e grupo de estudantes concluintes e ingressantes que participaram do ENADE/2009 na área de SECRETARIADO EXECUTIVO.
Tanto para concluintes quanto para ingressantes, a maioria considerou a extensão da prova como Adequada em relação ao tempo destinado à sua resolução, respectivamente 58,3 \% e 62,1 \%.
Entre os concluintes esta proporção cresceu com o desempenho.
Entre os ingressantes, esta proporção, também, é crescente.
A proporção dos que consideraram a extensão da prova muito curta foi ligeiramente maior entre os concluintes ( 0,6 \% ) do que entre os ingressantes ( 0,1 \% ).
A proporção dos que consideraram a extensão da prova curta/muito curta vis-à-vis ao tempo total foi maior entre os concluintes ( 3,2 \% ) do que entre os ingressantes ( 2,3 \% ).
Entre os ingressantes, somente os estudantes com desempenho no primeiro e segundo grupo declararam que a extensão da prova teria sido muito curta com relação ao tempo total.
A maior proporção de concluintes declarando que a prova foi curta/muito curta foi entre os estudantes de pior desempenho, 4,7 \%.
A proporção dos que consideraram a extensão da prova muito longa foi próxima entre os concluintes ( 12,0 \% ) e ingressantes ( 11,2 \% ).
A proporção dos que consideraram a extensão da prova longa foi menor entre os concluintes ( 26,6 \% ) do que entre os ingressantes ( 24,4 \% ).
Os Gráficos 4.6a e 4.6b apresentam a distribuição das respostas à questão sobre a extensão da prova em relação ao tempo total, segundo a Região onde o curso está situado para os estudantes concluintes e ingressantes que participaram do ENADE/2009 na área de SECRETARIADO EXECUTIVO.
Em todas as Regiões brasileiras, a maioria dos estudantes, tanto concluintes como ingressantes, perceberam a prova como Adequada.
Somente os concluintes da Região Sudeste e Sul ( 61,2 \% e 60,3 \%, respectivamente ) se situaram acima da média ( 58,3 \% ) para esta resposta, enquanto entre os ingressantes, os da Região Norte ( 67,3 \% ), Sudeste ( 65,0 \% ) e Sul ( 62,9 \% ) ficaram acima da média ( 62,1 \% ). 4.3 Compreensão dos Enunciados das Questões 4.3.1 Formação Geral Os Gráficos 4.7a, 4.7b e 4.7c apresentam a distribuição das respostas à questão “ Os enunciados das questões da prova na parte de Formação Geral, estavam claros e objetivos? ”, segundo quartos de nota final e grupo de estudantes concluintes e ingressantes que participaram do ENADE/2009 na área de SECRETARIADO EXECUTIVO.
Com relação aos enunciados das questões da parte de Formação Geral, as opiniões foram positivas : a grande maioria dos concluintes ( 77,2 \% ) e ingressantes ( 78,3 \% ) consideraram claros e objetivos os enunciados de todas ou da maioria das questões.
Entre os concluintes, esta proporção cresceu com o desempenho : 68,8 \% para o grupo com o pior desempenho e 81,9 \% para o de melhor.
Entre os ingressantes, esta proporção cresce de igual maneira, indo de 71,3 \% a 85,9 \%.
Os Gráficos 4.8a e 4.8b apresentam a distribuição das respostas à questão “ Os enunciados das questões da prova na parte de Formação Geral, estavam claros e objetivos? ”, segundo a Região onde o curso está situado para os estudantes concluintes e ingressantes que participaram do ENADE/2009 na área de SECRETARIADO EXECUTIVO.
Em termos regionais, as proporções de estudantes que consideraram claros e objetivos os enunciados de todas ou da maioria das questões variaram entre os concluintes, de 81,2 \% no Norte a 74,2 \% no Centro-Oeste.
Já entre os ingressantes estas proporções variaram de 81,5 \% no Norte a 76,3 \% no Sudeste.
Com relação às respostas “ Poucos ” e “ Não, nenhum ” entre concluintes as maiores proporções ocorreram nas Regiões Nordeste e Sul ( 9,6 \% ) e entre ingressantes na Região Centro-Oeste ( 10,0 \% ). 4.3.2 Componente Específico Os Gráficos 4.9a, 4.9b e 4.9c apresentam a distribuição das respostas à questão “ Os enunciados das questões da prova na parte de formação específica estavam claros e objetivos? ”, segundo quartos de nota final e grupo de estudantes concluintes e ingressantes que participaram do ENADE/2009 na área de SECRETARIADO EXECUTIVO.
Com relação aos enunciados das questões da parte do Componente Específico, as opiniões foram também positivas : a grande maioria dos concluintes ( 78,6 \% ) e ingressantes ( 77,0 \% ) consideraram claros e objetivos os enunciados de todas ou da maioria das questões.
A proporção de ingressantes declarando que todos os enunciados estavam claros e objetivos foi decrescente com o desempenho indo de 22,5 \% no primeiro quarto para 16,6 \% para aqueles com melhor desempenho.
O número de estudantes que declaram que “ Não, nenhum ” dos enunciados das questões da parte do Componente Específico era claro e objetivo, foi equivalente com a parte de Formação Geral : 0,5 \% entre os concluintes e 0,7 \% entre os ingressantes.
Os que responderam Poucos são 6,5 \% entre os concluintes e 7,9 \% entre os ingressantes.
Os Gráficos 4.10a e 4.10b apresentam a distribuição das respostas à questão “ Os enunciados das questões da prova na parte de formação específica estavam claros e objetivos? ”, segundo a Região onde o curso está situado para os estudantes concluintes e ingressantes que participaram do ENADE/2009 na área de SECRETARIADO EXECUTIVO.
Em termos regionais, as proporções de estudantes que consideraram claros e objetivos os enunciados de todas ou da maioria das questões variaram entre os concluintes, de 74,4 \% no Centro-Oeste a 81,3 \% no Nordeste.
Já entre os ingressantes estas proporções variaram de 74,1 \% na Região Sudeste a 84,3 \% no Norte e as Regiões Sul e Centro-Oeste, também, apresentaram valores acima da média.
Com relação às respostas “ Poucos ” e “ Não, nenhum ” para concluintes a maior proporção ocorreu na Região Centro-Oeste com 9,1 \%.
Para os ingressantes, as Regiões que apresentaram valores acima da média ( 8,5 \% ) foram a Sudeste e a Centro-Oeste, com valores respectivamente de 10,1 \% e 10,1 \%. 4.4 Suficiência das informações/instruções fornecidas Os Gráficos 4.11a, 4.11b e 4.11c apresentam a distribuição das respostas à questão “ As informações/instruções fornecidas para a resolução das questões foram suficientes para resolvê -las? ”, segundo quartos de nota final e grupo de estudantes concluintes e ingressantes que participaram do ENADE/2009 na área de SECRETARIADO EXECUTIVO.
Quando perguntados sobre as informações/instruções fornecidas para a resolução das questões, quase 85,5 \% dos concluintes e 82,5 \% dos ingressantes de todo o Brasil responderam que estas eram ou excessivas ou suficientes em todas ou na maioria das questões, demonstrando uma avaliação positiva da prova.
Levando -se em conta o desempenho dos participantes, percebeu -se consideráveis semelhanças entre as opiniões de estudantes, ingressantes e concluintes.
A proporção de estudantes com a opinião supracitada foi crescente com o desempenho, tanto para concluintes como para ingressantes.
Entre os concluintes a proporção vai de 79,9 \% para os de pior desempenho a 90,2 \% para os de melhor desempenho.
Entre os ingressantes as proporções equivalentes foram, respectivamente, 75,9 \% a 89,3 \%.
O complemento desta opinião, ou seja, os estudantes que responderam somente em algumas ou em nenhuma delas, apresentou uma proporção decrescente com o desempenho tanto para os concluintes como para os ingressantes.
A resposta “ não, em nenhuma delas ”, não apareceu apenas entre os ingressantes com o melhor desempenho.
Os Gráficos 4.12a e 4.12b apresentam a distribuição das respostas à questão “ As informações/instruções fornecidas para a resolução das questões foram suficientes para resolvê -las? ”, segundo a Região onde o curso está situado para os estudantes concluintes e ingressantes que participaram do ENADE/2009 na área de SECRETARIADO EXECUTIVO.
A Região Centro-Oeste foi a que apresentou os percentuais mais baixos, com 80,1 \% dos ingressantes e 82,2 \% dos concluintes emitindo tal opinião.
Os Gráficos 4.12a e 4.12b apresentam a distribuição das respostas à questão “ As informações/instruções fornecidas para a resolução das questões foram suficientes para resolvê -las? ”, segundo a Região onde o curso está situado para os estudantes concluintes e ingressantes que participaram do ENADE/2009 na área de SECRETARIADO EXECUTIVO.
A Região Centro-Oeste foi a que apresentou os percentuais mais baixos, com 80,1 \% dos ingressantes e 82,2 \% dos concluintes emitindo tal opinião. 4.5 Dificuldade encontrada ao responder à prova Os Gráficos 4.13a, 4.13b e 4.13c apresentam a distribuição das respostas à questão “ Você se deparou com alguma dificuldade ao responder a prova? ”, segundo quartos de nota final e grupo de estudantes concluintes e ingressantes que participaram do ENADE/2009 na área de SECRETARIADO EXECUTIVO.
Perguntados sobre as dificuldades com as quais se depararam ao responder à prova, 42,6 \% dos ingressantes apontaram como principal dificuldade o desconhecimento do conteúdo, já os 36,2 \% dos concluintes apontaram como principal dificuldade a forma diferente de abordagem.
Na verdade para a quase totalidade das áreas, ingressantes apresentaram uma proporção de estudantes bem maior do que entre os concluintes declarando como principal dificuldade o desconhecimento da matéria.
O desconhecimento do conteúdo foi apontado somente por 13,5 \% dos concluintes.
A forma diferente de abordagem do conteúdo foi considerada como maior dificuldade por 27,5 \% dos ingressantes, enquanto que a falta de motivação foi a principal dificuldade para 19,9 \% dos concluintes e para 11,6 \% dos ingressantes.
Entre os concluintes, esta dificuldade apresentou valores decrescentes conforme o desempenho.
Em termos nacionais, quase uma proporção maior de concluintes do que de ingressantes declararam que não tiveram qualquer dificuldade para responder à prova, respectivamente 13,2 \% e 8,2 \%.
Os Gráficos 4.14a e 4.14b apresentam a distribuição das respostas à questão “ Você se deparou com alguma dificuldade ao responder a prova? ”, segundo a Região onde o curso está situado para os estudantes concluintes e ingressantes que participaram do ENADE/2009 na área de SECRETARIADO EXECUTIVO.
Em relação às Regiões brasileiras, o percentual de ingressantes que apontaram o desconhecimento do conteúdo como principal dificuldade ao responder à prova variou de 37,4 \% na Região Centro-Oeste, a 54,5 \% na Sul.
Entre os concluintes, tal percentual variou de 7,9 \% na Região Norte a 16,2 \%, na Sul.
A forma diferente de abordagem de conteúdo foi citada como maior dificuldade por 29,7 \% dos concluintes na Região Nordeste chegando a 40,9 \% na Região Norte.
Já entre os ingressantes os valores foram de 21,2 \% na Região Sul a 33,2 \% na Região Centro-Oeste.
O percentual de concluintes que citaram a falta de motivação como principal dificuldade variou de 16,5 \% na Região Sudeste a 25,4 \% na Região Nordeste.
Entre os ingressantes as estatísticas correspondentes foram de 9,8 \% na Região Sul a 14,9 \%, na Região Centro-Oeste.
Os que declararam que não tiveram qualquer tipo de dificuldade variaram, entre os concluintes, de 11,6 \% na Região Norte a 14,9 \% na Região Sudeste.
Entre os ingressantes este percentual variou de 6,1 \% na Região Sul a 9,3 \% na Região Sudeste. 4.6 Percepção sobre os conteúdos das questões objetivas da prova Os Gráficos 4.15a, 4.15b e 4.15c apresentam a distribuição das respostas à questão “ Considerando apenas as questões objetivas da prova, você percebeu que :... ”, segundo quartos de nota final e grupo de estudantes concluintes e ingressantes que participaram do ENADE/2009 na área de SECRETARIADO EXECUTIVO.
As respostas possíveis para esta questão estavam relacionadas à cobertura do conteúdo das questões objetivas durante o curso.
A proporção de ingressantes ( 55,1 \% ) que declarou que não estudou ainda a maioria desses conteúdos foi bem maior do que entre os concluintes ( 4,9 \% ).
Esta proporção foi decrescente tanto para os concluintes como para os ingressantes como função do desempenho.
Por outro lado, a proporção de estudantes que declarou que estudou e aprendeu todos ou a maioria desses conteúdos, entre os ingressantes foi de 25,3 \% e entre os concluintes 67,1 \%.
Esta proporção foi crescente como função do desempenho, tanto para os concluintes como para os ingressantes.
A proporção de estudantes que declarou ter estudado a maioria desses conteúdos, mas não os ter aprendido, foi bem maior entre os concluintes ( 16,5 \% ) do que os ingressantes ( 4,8 \% ).
Os Gráficos 4.16a e 4.16b apresentam a distribuição das respostas à questão “ Considerando apenas as questões objetivas da prova, você percebeu que :... ”, segundo a Região de habilitação do curso para os estudantes concluintes e ingressantes que participaram do ENADE/2009 na área de SECRETARIADO EXECUTIVO, segundo a Região onde o curso está situado.
Em relação às Regiões brasileiras, concluintes que afirmaram não ter estudado ainda a maioria dos conteúdos ocorreram em todas as Regiões, os percentuais variaram de 3,5 \% na Região Nordeste a 6,9 \%, na Centro-Oeste.
Entre os ingressantes, os percentuais variaram de 45,2 \% na Região Sudeste a 70,5 \%, na Norte.
Em todas as Regiões, exceto a Região Norte, a maioria dos concluintes afirmou ter estudado e aprendido todos ou muitos dos conteúdos, com percentuais que variaram de 59,5 \% na Região Centro-Oeste a 71,5 \%, na Sudeste.
Para os ingressantes, tal percentual variou de 15,2 \%, na Região Norte, a 33,3 \%, na Sudeste. 4.7 Tempo gasto para concluir a prova Os Gráficos 4.17a, 4.17b e 4.17c apresentam a distribuição das respostas à questão “ Qual foi o tempo gasto para concluir a prova? ”, segundo quartos de nota final e grupo de estudantes concluintes e ingressantes que participaram do ENADE/2009 na área de SECRETARIADO EXECUTIVO.
Quando perguntados quanto ao tempo de conclusão da prova, a maioria respondeu entre duas e três horas, tanto entre os concluintes ( 39,9 \% ) quanto entre os ingressantes ( 39,1 \% ).
A segunda maior incidência de respostas foi entre uma e duas horas, igualmente para ambos os grupos : concluintes ( 30,1 \% ) e ingressantes ( 33,2 \% ).
A proporção daqueles que declararam ter usado as quatro horas mas não terem conseguido terminar foi bem baixo : 2,0 \% entre os ingressantes e 2,6 \% entre os concluintes.
Os Gráficos 4.18a e 4.18b apresentam a distribuição das respostas à questão “ Qual foi o tempo gasto para concluir a prova? ”, segundo a Região onde o curso está situado para os estudantes concluintes e ingressantes que participaram do ENADE/2009 na área de SECRETARIADO EXECUTIVO.
A proporção de concluintes que gastaram entre uma e duas horas foi maior na Região Sul ( 36,2 \% ) e menor na Região Centro-Oeste ( 22,3 \% ).
Para ingressantes, o maior percentual para este tempo ocorreu na Região Centro-Oeste ( 42,8 \% ) e a menor na Região Norte ( 17,8 \% ).

% matched lemmas: 
\end{textsample}
